% #############################################################################
% This is Chapter 1 - Introduction
% !TEX root = ../main.tex
% #############################################################################
% Written as one continuous argument with no sections, to hold the chapter to
% two pages. sec:motivation and sec:intro_objectives are kept as inline labels
% because Chapters 3, 6, 9 and 10 reference them; they resolve to this chapter.
% #############################################################################
\fancychapter{Introduction}
\label{chap:intro}

\ac{GenAI} and \acp{LLM} have extended machine assistance towards generating coherent, context-aware artefacts, introducing intent-driven programming, colloquially \emph{Vibe Coding}, in which developers describe desired behaviour in natural language and refine what a model produces~\cite{Moore:2025vc, Zamfiroiu:2025ec}. The shift is not confined to code: \acp{LLM} now participate in requirements elicitation, architectural exploration, test generation, documentation and debugging, touching every \ac{SDLC} phase~\cite{Ahmed:2025ai, Zhang:2025ea, Hou:2025cl}. The practitioner's work moves from constructing software towards specifying intent, supervising generated artefacts and remaining accountable for them.

\label{sec:motivation}%
The evidence on outcomes, however, is more nuanced than early enthusiasm suggested, and that gap motivates this work. Adoption is close to universal while trust is not: 84\% of developers report using or planning to use AI tools, up from 76\%, while the proportion trusting their accuracy fell from 40\% to 29\%~\cite{StackOverflow:2025dev}, with daily use around 90\%~\cite{DORA:2025ai}. That distrust is warranted, and the decisive point is that functional correctness and security do not move together: 61\% of a leading coding agent's solutions were functionally correct but only 10.5\% secure~\cite{Zhao:2025vibe}, and across more than one hundred models 45\% of samples introduced an OWASP Top 10 vulnerability, with newer models improving at correctness while security stayed flat~\cite{Veracode:2025gen}.

Productivity gains are equally conditional. Early evidence reported a task completed 55.8\% faster with an assistant, though unpublished in a peer-reviewed venue~\cite{Peng:2023cop}, whereas a randomised controlled trial found 16 experienced developers on their own mature repositories took 19\% longer while estimating afterwards that they had been 20\% faster~\cite{Becker:2025rct}; the favourable result came from a self-contained task with no codebase to respect. At organisational scale, telemetry across 1,255 teams associates high adoption with 98\% more pull requests merged but 91\% longer review times and no significant correlation with delivery outcomes~\cite{Faros:2025par}. The bottleneck moves downstream to validation rather than disappearing. Several sources here are industry-reported rather than peer-reviewed and are identified as such throughout, so the argument rests on consistency of direction rather than any single figure: the value of AI assistance is conditional on the process surrounding it rather than a property of the model alone.

To characterise practice rigorously, an \ac{SLR} following Kitchenham's guidelines retrieved 646 articles, reported in \Cref{chap:slr}. Its central finding is methodological: organisations lack established methodologies for the responsible use of \ac{AI} in professional software development~\cite{Banh:2025cf, Farhanna:2025eu, Hemmat:2025rd, Gao:2025cc}, and proposed approaches remain fragmented across isolated tools, sharing no common foundation on lifecycle integration, oversight or evaluation~\cite{Ahmed:2025ai}. Established methodologies do not close the gap: Waterfall, Scrum, Kanban and the Agile Manifesto were designed for environments in which every contributor is human~\cite{Zhang:2025ea}, so they lack structures for governing AI participation, bounding the context a model may act upon, and allocating responsibility over generated artefacts, while their cadence assumes implementation is the binding constraint, which the evidence contradicts. \textbf{The research problem is therefore the absence of structured, repeatable guidance and standardised evaluation practices for integrating \ac{AI} across the \ac{SDLC} while preserving human oversight, software quality and security, and accountability.}

\label{sec:intro_objectives}%
The objective is to design, instantiate, demonstrate and evaluate a process-oriented framework addressing that problem. Four questions structure it, labelled \emph{\ac{DRQ}} to keep them distinct from the review questions of \Cref{chap:slr}, which ask what the literature reports rather than what these artefacts achieve. \textbf{DRQ1:} how can AI-supported activities be mapped to \ac{SDLC} phases so that the point, purpose and boundary of AI participation are unambiguous? \textbf{DRQ2:} which governance mechanisms, quality gates and human-in-the-loop practices keep AI contributions traceable, validated and accountable? \textbf{DRQ3:} to what extent can those constructs be operationalised in a working system, and which resist implementation? \textbf{DRQ4:} does the framework, through that implementation, give practitioners usable and cognitively sustainable support in a real setting? \Cref{chap:proposal,chap:apex,chap:demonstration,chap:evaluation} answer them in turn, the third including the constructs weakened, proxied or left unimplemented.

The contributions are four: the synthesis above; a process framework treating AI as a first-class collaborator, with phases, gates, function-based responsibilities, metrics and a versioned specification anchoring artefacts to requirements; Apex, a reference implementation showing those constructs implementable rather than describable; and an empirical evaluation of its usability and workload.
